\chapter{S'entraîner avec \texttt{forge-apprendre}}\label{ch:apprendre}

\begin{objectifs}[Objectifs]
\begin{itemize}
  \item suivre quatre leçons pratiques sur un cas fictif, avec un coach ;
  \item mesurer votre progression sans fichier de suivi ;
  \item transférer chaque leçon sur votre projet réel.
\end{itemize}
\end{objectifs}

\begin{situation}[Sur le terrain : une nouvelle recrue]
Une ingénieure rejoint l'équipe de \sika{}. Elle ne veut pas lire un manuel de bout en bout : elle veut essayer. Le skill
\skill{forge-apprendre} lui donne un coach et un terrain d'essai.
\end{situation}

\section{Un coach, pas un exécutant}
Le skill \emph{explique}, \emph{relit} et donne des indices à la demande ; c'est \textbf{vous} qui écrivez chaque
livrable. Les exercices portent sur un autre cas fictif que \sika{}, \textbf{PayBoutique} (une boutique en ligne béninoise
qui encaisse des paiements MTN MoMo et Moov Money), et vivent dans \cmd{apprentissage/}, jamais dans \cmd{docs/} : le vrai dossier
de conception reste intact.

\begin{concept}[Apprentissage par la pratique]
Chaque leçon suit le même schéma : le \emph{concept} en quelques lignes, le \emph{pourquoi}, puis un \emph{exercice} vérifié par un garde-fou. À la fin,
vous refaites l'exercice sur votre projet avec le skill de l'étape.
\end{concept}

\section{Les quatre leçons}
\noindent\begin{tabularx}{\linewidth}{lXX}
\toprule
\textbf{Leçon} & \textbf{Exercice} & \textbf{Vérification} \\ \midrule
1 Impact map & objectif mesurable, acteur, tri de trois idées & mots-clés \emph{indicateur}, \emph{cible}, \emph{échéance} \\
2 Langage & glossaire et code contenant un terme banni & \cmd{check-glossary-terms.sh} : KO puis OK \\
3 Spec & règle de paiement avec trois cas & \cmd{check-spec-validated.sh} \\
4 ADR & choix d'hébergement chiffré en FCFA & \cmd{check-adr-present.sh} \\
\bottomrule
\end{tabularx}

\section{Suivre sa progression}
\cmd{scripts/forge-lecons.sh} affiche par exemple \cmd{Leçons [██░░] 2/4 · prochaine : 03-spec}, calculé à partir de vos fichiers d'exercice : vous
reprenez là où vous vous étiez arrêté, sans fichier de suivi.

\begin{piege}[Piège : croire que le vert prouve la qualité]
Les scripts ne vérifient que la présence et le statut, pas la qualité. La qualité se discute en relecture avec le coach. Les concepts
sans leçon (bounded context, OpenAPI d'abord, fitness functions, recette) s'étudient à partir du skill correspondant.
\end{piege}

\begin{vousjouez}[À vous de jouer]
Lancez \skill{forge-apprendre}, faites la leçon~1, puis refaites immédiatement l'exercice sur votre projet avec \skill{forge-cadrage}.
Qu'est-ce qui était différent quand les chiffres étaient les vôtres ?
\end{vousjouez}

\begin{retenir}[À retenir]
\begin{itemize}
  \item Le coach explique ; vous écrivez.
  \item Quatre leçons : impact map, langage, spec, ADR.
  \item Après chaque leçon, transférez sur votre projet avec le skill de l'étape.
\end{itemize}
\end{retenir}
\source{skills/forge-apprendre/SKILL.md et lecons/}

\begin{acquis}
  \item Où vivent les exercices, et pourquoi pas dans \cmd{docs/} ?
  \item Que ne vérifie \emph{pas} \cmd{forge-lecons.sh} ?
  \item Quelle leçon se rattache à quel skill d'étape ?
\end{acquis}
